\documentclass[twocolumn]{aastex631}
\begin{document}
\title{The reionization ionizing-photon-budget shortfall for star-forming galaxies is not robust to systematics at z\~{}6 (lowxi systematic corner)}
\author{NebulaMind Lab (autonomous pipeline)}
\affiliation{NebulaMind Lab, \url{https://lab.nebulamind.net}}
\begin{abstract}
Reionization ionizing-photon-budget reconciliation at z\~{}6: star-forming galaxies require f\_esc=0.062 (+0.062/-0.032) to close the budget (Madau-Dickinson SFRD, log xi\_ion=25.5+/-0.15, clumping C=2-5, JWST-SFRD tail), versus indirect-proxy-inferred f\_esc=0.062 (+0.108/-0.039) from LzLCS O32/beta calibrations. Median delta(required-inferred)=-0.001 dex-frac (16-84\%: -0.107 to +0.067); 50\% of the systematic MC shows a shortfall. Conclusion: the budget CLOSES within the systematic, and the sign flips sign between the O32 and beta calibrations (calibration-driven, not robust). The result is bounded by the xi\_ion x clumping x proxy-calibration systematic, not by statistics. Generated autonomously from public data (jwst) via the NebulaMind Lab runner: a bounded, reproducible, descriptive study.
\end{abstract}
\section{Introduction}
Recent studies have highlighted potential discrepancies in the ionizing photon budget during reionization, raising concerns about whether star-forming galaxies can account for the necessary photons to drive this process [Muñoz2024]. This issue is further complicated by uncertainties in key parameters such as the escape fraction of ionizing photons (f\_esc) and the clumping factor of ionized hydrogen (C). Previous works have explored various aspects of reionization, including the role of absorption-dominated scenarios [Davies2021] and the calibration of excursion set models [Park2022]. However, a comprehensive analysis of the photon budget is still needed to reconcile these discrepancies.
\section{Data and method}
To address this issue, we employ a literature-anchored approach that leverages established values from previous research. Specifically, we use the cosmic star formation rate density (SFRD) from Madau \& Dickinson (2014), along with published calibrations for xi\_ion and O32/beta f\_esc proxy [LzLCS: Chisholm+22, Flury+22; Simmonds+24]. By combining these elements, we calculate the ionizing photon budget at z\~{}6 using a systematic method that avoids reliance on specific survey catalog data.
\section{Result}
Our analysis reveals that star-forming galaxies can close the reionization ionizing-photon-budget at z\~{}6 if they have an escape fraction of f\_esc=0.062 (+0.062/-0.032). This value is consistent with indirect-proxy-inferred estimates from LzLCS O32/beta calibrations, which yield f\_esc=0.062 (+0.108/-0.039). The median difference between the required and inferred escape fractions is -0.001 dex-frac, ranging from -0.107 to +0.067 (16-84\% confidence interval), indicating that 50\% of systematic Monte Carlo simulations show a shortfall.
\begin{figure}\centering\includegraphics[width=\columnwidth]{result.png}\caption{The reionization ionizing-photon-budget shortfall for star-forming galaxies is not robust to systematics at z\~{}6 (lowxi systematic corner)}\end{figure}
\section{Caveats}
Despite these findings, it is essential to acknowledge the limitations of our approach. Our analysis relies on automated, single-selection measurements and lacks calibration, which may introduce biases or uncertainties in the results. Furthermore, the ionizing photon budget is sensitive to various systematics, including the assumed values for xi\_ion, clumping factor C, and proxy-calibration choices. These factors can significantly impact our understanding of reionization, emphasizing the need for further research and refinement of these parameters [Madau2017].
\section*{References}
\footnotesize
[Muoz2024] Muñoz et al. (2024). Reionization after JWST: a photon budget crisis?. bibcode 2024MNRAS.535L..37M \par [Davies2021] Davies et al. (2021). The Predicament of Absorption-dominated Reionization: Increased Demands on Ionizing Sources. bibcode 2021ApJ...918L..35D \par [Park2022] Park et al. (2022). Calibrating excursion set reionization models to approximately conserve ionizing photons. bibcode 2022MNRAS.517..192P \par [Duncan2015] Duncan et al. (2015). Powering reionization: assessing the galaxy ionizing photon budget at z < 10. bibcode 2015MNRAS.451.2030D \par [Madau2017] Madau (2017). Cosmic Reionization after Planck and before JWST: An Analytic Approach. bibcode 2017ApJ...851...50M

\end{document}
